\section{Average Number of Solutions}\label{sec:average-solutions}
We now investigate the \emph{average number of solutions}\index{average number of solutions} of a polynomial equations. Let $d\in\bbN$. Then define $\Om_d\coloneqq \bbF_q^{\leq d}[X_1,\dots, X_n]$ and $\om(d)$ be the set of $n-$tuples $(i_1,\dots,i_n)\in\bbZ_0$ such that $i_1+\cdots +i_n\leq d$. Then 
\begin{observation}\label{obs:Omd-omd}
  With the definitions of $\Om_d$ and $\om(d)$ as above we have $|\Om_d|=q^{|\om(d)|}$. 
\end{observation}
\begin{lemma}{Average of Solution Count}{lem:average-solutions}
  \index{average number of solutions}%
  With the notations of $\Om_d$ and $\om(d)$ as defined above we have $$\frac1{|\Om_d|}\sum_{f\in\Om_d}Z(f)=q^{n-1}$$
\end{lemma}
\begin{proof}
  We have $$\sum_{f\in\Om_d}Z(f)=\sum_{f\in\Om_d}\sum_{\substack{\alpha_1,\dots,\alpha_n\in\bbF_q\\ f(\alpha_1,\dots,\alpha_n)=0}}1=\sum_{\alpha_1,\dots,\alpha_n\in\bbF_q}\sum_{\substack{f\in\Om_d\\ f(\alpha_1,\dots,\alpha_n)=0}}1$$Let for any $f\in\Om_d$ $$f(X_1,\dots,X_n)=\sum_{(i_1,\dots,i_n)\in\om(d)}f_{i_1,\dots,i_n}X_1^{i_1}\cdots X_n^{i_n}\quad f_{i_1,\dots,i_n}\in\bbF_q,\ \forall\ (i_1,\dots, i_n)\in\om(d)$$Then for a fixed $\alpha_1,\dots, \alpha_n\in\bbF_q$ we get to choose the coefficients $f_{i_1,\dots, i_n}$ for all $(i_1,\dots,i_n)\in\om(d)\setminus\{(0,\dots, 0)\}$ arbitrarily and then determining the constant term $f_{0,\dots, o}$ to make $f(\alpha_1,\dots,\alpha_n)=0$. Thus the number of $f\in\om_d$ with $f(\alpha_1,\dots, \alpha_n)=0$ is equal to $q^{|\om(d)|-1}$. Therefore $\sum_{f\in\Om_d}Z(f)=q^n\cdot q^{|\om(d)|-1}=|\Om_d|\cdot q^{n-1}$. So we have the lemma. 
\end{proof}
The proof of the above lemma gives us the following observation:
\begin{observation}\label{obs:num-poly-zero-at-point}
  For any $\alpha_1,\dots,\alpha_n\in\bbF_q$, the number of polynomials $f\in\Om_d$ such that $f(\alpha_1,\dots,\alpha_n)=0$ is $q^{|\om(d)|-1}$. 
\end{observation}

The above lemma implies that a polynomial equation in $n$ variables has on average $q^{n-1}$ solutions in $\bbF_q^n$. So we now consider average deviation from the average number of solutions, $q^{n-1}$. 
\begin{Theorem}{Variance of Solution Count}{}
  \index{solution count!variance}%
  With the notations of $\Om_d$ and $\om(d)$ as defined above we have
  $$\frac1{|\Om_d|}\sum_{f\in\Om_d}\Big(Z(f)-q^{n-1}\Big)^2=q^{n-1}-q^{n-2}$$
\end{Theorem}
\begin{proof}
  Let $\alpha=(\alpha_1,\dots,\alpha_n),\beta=(\beta_1,\dots,\beta_n)\in\bbF_q^n$ be any two points. Then we get
  $$\sum_{f\in\Om_d}Z(f)^2  = \sum_{f\in\Om_d}\lt( \sum_{\substack{\alpha\in\bbF_q^n\\ f(\alpha)=0}} 1 \rt)^2= \sum_{f\in\Om_d}\sum_{\substack{\alpha\in\bbF_q^n\\ f(\alpha)=0}}\sum_{\substack{\beta\in\bbF_q^n\\ f(\beta)=0}}1=\sum_{\alpha,\beta\in\bbF_q^n}\sum_{\substack{f\in\Om_d\\  f(\alpha)=f(\beta)=0}}1$$
  Now if $\alpha=\beta$ then by \obsref{obs:num-poly-zero-at-point} we have $\sum_{f\in\Om_d,\   f(\alpha)=0}1=q^{|\om(d)|-1}$. Now suppose $\alpha\neq\beta$. This gives two linear equations for the coefficients of $f$. Therefore there are $q^{|\om(d)|-2}$ polynomials $f\in\Om_d$ such that $f(\alpha)=f(\beta)=0$. Hence we obtain \begin{align*}
    \sum_{f\in\Om_d}Z(f)^2  =  & = \sum_{\alpha\in\bbF_q^n}|\Om_d|\cdot q^{-1}+\sum_{\alpha,\beta\in\bbF_q^n,\alpha\neq \beta}|\Om_d|\cdot q^{-2} \byo{Omd-omd}\\ 
    & = q^n\cdot |\Om_d|\cdot q^{-1}+q^n(q^n-1)|\Om_d|\cdot q^{-2}\\ 
    & = |\Om_d|\Big( q^{n-1}+q^{2n-2}-q^{n-2} \Big)
  \end{align*}
  So now we get \begin{align*}
    \sum_{f\in\Om_d}(Z(f)-q^{n-1})^2 & = \sum_{f\in\Om_d}Z(f)^2 + 2\cdot q^{n-1}\sum_{f\in\Om_d}Z(f)+q^{2n-2}\sum_{f\in\Om_d}1\\ 
    & = |\Om_d|\Big(q^{2n-2}+q^{n-1}-q^{n-2}\Big)-2\cdot q^{n-1}\cdot |\Om_d|\cdot q^{n-1}+q^{2n-2}\cdot |\Om_d|\\ 
    & = |\Om_d|(q^{n-1}-q^{n-2})
  \end{align*}
  So we have the result. 
\end{proof}
From this result we can expect that $|Z(f)-q^{n-1}|$ is often $O\lt(q^{\nicefrac{(q-1)}{2}}\rt)$. Later we will see instances of such expected behavior for various polynomial equations. In the next sections we will talk about character sums over polynomial evaluations. 
